% Diapositiva 22
% David

%=============================================
% 6. Explicar que el mínimo representa la mejor estimación
%=============================================
\begin{frame}{Optimización: Interpretación del Mínimo Global}

    \begin{columns}[T]
        %========================================
        % COLUMNA IZQUIERDA: El Objetivo Matemático
        %========================================
        \column{0.48\textwidth}
        \begin{block}{Criterio de Optimización}
            \scriptsize
            Determinación del punto óptimo en el espacio tridimensional mediante la minimización del error global:
            \vspace{0.3cm}
            \[
            (x^*, y^*, z^*) = \arg\min E(x,y,z)
            \]
            \vspace{0.2cm}
        \end{block}

        \vspace{0.2cm}

        \begin{block}{Validación del Modelo}
            \scriptsize
            Un valor residual bajo de $E(x,y,z)$ certifica que las coordenadas evaluadas reproducen con alta fidelidad las amplitudes medidas.
        \end{block}

        %========================================
        % COLUMNA DERECHA: ¿Qué significa realmente?
        %========================================
        \column{0.48\textwidth}
        \begin{exampleblock}{Significado Físico-Matemático}
            \scriptsize
            El punto mínimo de la función representa:
            \begin{itemize}
                \item \textbf{Ajuste Óptimo:} La menor discrepancia físicamente alcanzable respecto a la red de sensores.
                \vspace{0.3cm}
                \item \textbf{Consistencia:} Región espacial donde las predicciones teóricas convergen con los datos de campo.
            \end{itemize}
        \end{exampleblock}
    \end{columns}

    \vspace{0.5cm}

    \begin{center}
        \footnotesize
        \textbf{"   El mínimo no es solo un valor numérico decreciente; es la posición espacial que mejor explica el evento sísmico real."}
    \end{center}

\end{frame}